% !TeX root = RJwrapper.tex
\title{Working with table environments in texor}
\author{by Abhishek Ulayil}

\maketitle

\abstract{
This is a small sample article to demonstrate usage of \CRANpkg{texor} to convert table environments.
}

\section{Introduction}
Tables are commonly used in RJ articles to display data in a tabular format. However, there are differences in the way tables are handled by LaTeX and HTML.
LaTeX tables have more customization and are usually optimized for printing, whereas the web articles need tables optimized for varying sizes of media.
Pandoc converts most of the tables somewhat easily, but is unable to do well with table customization packages and complex tables.
Some pandoc extensions are used in order to tackle them, they are : \verb|simple_tables|,\verb|pipe_tables|
\section{Generic tables}
Normal LaTeX tables will be converted just fine, even if they include a bit of math and other elements.
However do note that the conversion will be done to traditional or pandoc style markdown instead of Rmarkdown
specific using \verb|knitr::kable(..)| function. Also any custom LaTeX commands or specific font charactars would
not be supported.
\begin{table}[htbp]
\centering
\begin{tabular}{l | llll }
 \hline
 Graphics Format & LaTeX & Markdown & Rmarkdown & HTML \\
 \hline
 PNG       & Yes & Yes & Yes & Yes \\
 JPG       & Yes & Yes & Yes & Yes \\
 PDF       & Yes & No & No & No \\
 SVG       & No & Yes & Yes & Yes \\
 Tikz      & Yes & No & Yes & No \\
 Algorithm & Yes & No & No & No \\
\hline
\end{tabular}
\caption{Image Format support in various Markup/Typesetting Languages}
\label{table:1}
\end{table}

\section{Multicolumn tables}
\begin{table*}[htbp]
\begin{center}
\begin{tabular}{l | llll }
 \hline
 \multicolumn{1}{c |}{EXAMPLE} & \multicolumn{2}{c}{$X$} &
\multicolumn{2}{c}{$Y$} \\
 \hline
  & 1 & 2 & 1 & 2 \\
 EX1  & X11 & X12 &  Y11  & Y12 \\
 EX2  & X21 & X22 &  Y21  & Y22 \\
 EX3  & X31 & X32 &  Y31  & Y32 \\
 EX4  & X41 & X42 &  Y41  & Y42\\
 EX5  & X51 & X52 &  Y51  & Y52 \\
\hline
\end{tabular}
\caption{An Example Table with Multicolumn}
\label{table:2}
\end{center}
\end{table*}

The above example illustrates a multicolumn table \ref{table:2}, which the \CRANpkg{texor} package can handle through subroutines and pre-processing step to transform the LaTeX code. 

\begin{itemize}
    \item The Stream Editor modifies the \verb|\multicolumn{..}| to \verb|\multicolumnx{..}|
    \item Then a LaTeX macro redefines the \verb|\multicolumnx{..}| to \verb|\multicolumn{---}|, which is accepted by pandoc.
    \item Finally pandoc reads the table and transforms it to markdown.
\end{itemize}

Also note that the stream editor is used to rename \verb|table*| environment to \verb|table| environment because the markdown/HTML will not support the specific changes table* environment will bring.

\section{Complex tables}
A complex table with various other elements like figures, math , code and so on, are also supported
by \CRANpkg{texor}.  
\begin{table}[htbp]
\centering
\begin{tabular}{|l | lll |}
 \hline
 Inline Format & LaTeX Support & Web Support &  Rendering \\
 \hline
 Text       & Yes & Yes  & \emph{Hello} \\ \hline
 Image       & Yes & Yes & \begin{minipage}{0.25\textwidth}
    \centering
    \includegraphics[width=1\textwidth]{Rlogo-5.png}
    \end{minipage} \\ \hline
 CodeBlock      & Yes & Yes & \begin{minipage}{0.45\textwidth}
\vspace{1mm}
\begin{example}
x <- 1:100
y <- dbinom(x,100,prob = 0.5)
plot(x,y)
\end{example}
    \end{minipage} \\ \hline
 Math      & Yes & Yes & $e = m c^2 $ \\ \hline
 Link      & Yes & Yes &  \href{www.google.com}{Google} \\ \hline
 Nested Table & Yes & No & NaN \\
\hline
\end{tabular}
\caption{Image Format support in various Markup/Typesetting Languages}
\label{table:3}
\end{table}

\section{Wide tables}

\begin{widetable}
\centering
\begin{tabular}{lllllllll}
\toprule
index  &  A  &  B     & C   &  D  &  E  &  F  &  G  & H  \\
\midrule
1 & 359.00 & NaN & 5796.00 & 0.00 & 16.14 & 1.00 & NaN & 0.00\\
2 & 25.73 & 0.00 & 1029.20 & NaN & 40.00 & 0.68 & 0.00 & 0.00\\
2.1 & 26.26 & 0.00 & 13.40 & 0.00 & 2.14 & 0.68 & 0.00 & NaN\\
2.2 & 32.06 & 20.06 & 47.64 & 0.04 & 1.80 & 0.68 & 0.01 & NaN\\
2.3 & 51.94 & 420.27 & 21.17 & 0.20 & 1.77 & 0.74 & 0.05 & NaN\\
2.4 & 40.62 & 30.44 & \phantom{000}0.90 & 0.57 & 1.44 & 1.31 & 0.24 & NaN\\
 & \ldots & \ldots&  \ldots& \ldots &  \ldots& \ldots& \ldots & \ldots \\
\bottomrule
\end{tabular}
\begin{tabular}{llllllll}
\toprule
index &    I  &  J  &  K  &  L  &  M  &  N & O \\
\midrule
1 & 0.00 & $-$1.73 & 0.00 & 0.00 & 0.00 & 0.00 & NaN\\
2 & 0.21 & $-$33.41 & 0.00 & 0.11 & 0.00 & NaN & $-$197.85\\
2.1 & 0.24 & $-$24.00 & 0.00 & 0.15 & NaN & 0.06 & $-$70.46\\
2.2 & 0.6 & $-$19.42 & 0.00 & 0.15 & 0.00 & 0.11 & $-$16.48\\
2.3 & 0.75 & $-$ 31.77 & 0.00 & 0.18 & 0.01 & 0.37 & $-$0.82\\
2.4 & 0.26 & $-$1.89 & 0.1 & 0.55 & 0.70 & 0.22 & $-$6.55 \\
  &\ldots &  \ldots&  \ldots& \ldots & \ldots & \ldots & \ldots \\
\bottomrule
\end{tabular}
\caption{\code{A dummy research data}}
\label{table:4}

\end{widetable}

Widetable environments also gets some perks when converted by \CRANpkg{texor} package. 
They get their own block and separate numbering (although references are shared by figures).

\begin{widetable}
\centering %
\begin{tabular}{llrr}
\toprule 
Package  & Commits  & Version  & Last Updated \tabularnewline
\toprule 
\CRANpkg{texor} & 260 & 1.1.0  & 28-Jul-2023 \tabularnewline
\CRANpkg{rebib} & 76 & 0.2.4  & 29-Jul-2023 \tabularnewline
\CRANpkg{rjtools}  & 314 & 1.0.11  & 30-Jul-2023 \tabularnewline
\CRANpkg{rmarkdown} & 3189 & 2.23  & 31-Jul-2023 \tabularnewline
\toprule 
\end{tabular}
\caption{A dummy summary of a few CRAN packages}
\label{table:5}

\end{widetable}


Here is a reference to the table \ref{table:4} and table \ref{table:5}.
\section{Summary}

In summary the \CRANpkg{texor} package supports:
\begin{itemize}
\item some common table environments.
\item wide tables.
\item some multicolumn tables.
\item inclusion of different environments like figure/code in tables.
\end{itemize}



\begin{thebibliography}{2}
    \providecommand{\natexlab}[1]{#1}
    \providecommand{\url}[1]{\texttt{#1}}
    \expandafter\ifx\csname urlstyle\endcsname\relax
      \providecommand{\doi}[1]{doi: #1}\else
      \providecommand{\doi}{doi: \begingroup \urlstyle{rm}\Url}\fi

\bibitem[Krewinkel, Lucero (2023)]{pandoc}
A.~ Krewinkel and A.~ Lucero
\newblock pandoc 3.0 Release notes
\newblock \emph{pandoc} \penalty0 2023
\newblock URL \url{https://pandoc.org/releases.html}

\end{thebibliography}


\address{%
Abhishek Ulayil\\
Student, Institute of Actuaries of India\\%
Mumbai, India\\
ORCiD: 0009-0000-6935-8690\\
}
